\documentclass{article}
\usepackage[a5paper,hmargin=1cm,vmargin=1.5cm]{geometry}
\usepackage{texmaths-symbols}
\usepackage{texmaths-picto}
\usepackage[ComicSansMS]{texmaths-fonts}

\long\def\writeword#1#2{%
    \noindent{\bfseries#1 :}\smallskip\par\noindent\begingroup
    \leftskip=20pt\rightskip=0pt plus 1fil\relax #2\par\endgroup
    \bigbreak
    }

\begin{document}
    Ce glossaire a pour but de dresser la liste des verbes utilisés dans les consignes
    des exercices mathématiques. L'élève qui se sentira démuni face à une consigne
    devra avoir comme premier réflexe de rechercher dans ce glossaire, le verbe de
    la consigne qui lui pose problème. Il y trouvera la définition.
    Afin d'ancrer ce glossaire dans la réalité de la classe, chaque verbe est
    accompagné d'une définition~:
    \begin{itemize}
        \item \guillemets{élève écrite} : celle que l'élève lira dans son glossaire~;
        \item \guillemets{orale élève} : celle reformulée à l'oral par l'enseignant si l'élève ne
        comprend pas la définition \guillemets{écrite élève}.
        \item  \guillemets{rigoureuse}: réservée à l'adulte (pas toujours nécessaire)~;
    \end{itemize}

    Le caractère absolu de ces définitions contredit la diversité des apprenants et
    des enseignants.

    Par conséquent, il n'est pas exclu de proposer, pour chacunes des trois catégories
    de définition, plusieurs définitions possibles. Cela permettra de s'adapter aux
    élèves/enseignants concernés.

    Les verbes seront classés par ordre alphabétique et il sera possible d'indiquer
    que la définition d'un verbe a été vue en classe entière (surlignage, soulignage,~...).
    L'élève n'aura qu'à chercher dans les verbes mis ainsi en valeur et non dans tout
    le glossaire, ce qui facilitera la recherche. On peut noter que l'élève s'exerce
    par la même occasion à effectuer une recherche par ordre alphabétique.

    Ce glossaire doit-il se restreindre aux seuls verbes ? Je pense qu'il pourrait
    s'étendre aux définitions \guillemets{implicites}, c'est à dire qui ne sont
    pas excplicitement écrites dans le cahier de leçon~? (par exemple, définition
    d'un outil géométrique).

    On peut penser à éventuellement laisser une partie ou (l'entièreté selon les élèves)
    de ce document (le document élève) en évaluations pour les aider.
    
    \newpage
    \begin{center}\LARGE
        Glossaire des verbes de consignes de mathématiques
    \end{center}
    \noindent\textit{Les verbes sont classés par ordre alphabétique.}
    \par\vskip10pt
    \input{glossaire-generated.tex}

\end{document}